% Generated draft: requires mathematical and source review.
\documentclass[11pt]{article}
\usepackage[margin=25mm]{geometry}
\usepackage{fontspec}
\setmainfont{Latin Modern Roman}
\usepackage{amsmath,amssymb,mathtools,unicode-math}
\setmathfont{Latin Modern Math}
\usepackage{enumitem,needspace,xurl,hyperref}
\setlength{\parindent}{0pt}
\setlength{\parskip}{5pt}
\setlength{\emergencystretch}{3em}
\newcommand{\N}{\mathbb{N}}
\newcommand{\Z}{\mathbb{Z}}
\newcommand{\Q}{\mathbb{Q}}
\newcommand{\R}{\mathbb{R}}
\newcommand{\C}{\mathbb{C}}
\newcommand{\F}{\mathbb{F}}
\newcommand{\A}{\mathbb{A}}
\newcommand{\PP}{\mathbb{P}}
\newcommand{\E}{\mathbb{E}}
\newcommand{\eps}{\varepsilon}
\newcommand{\dd}{\,\mathrm{d}}
\begin{document}
\section*{UnsolvedMath ID: 700647}
\textbf{Source problem code:} catalogue entry
\subsection{Definitions and mathematical statement}
For a group \(G\) and a group word \(w\) in \(k\) variables, its word map is
\[
 w_G:G^k\longrightarrow G,\qquad (g_1,\ldots,g_k)\longmapsto w(g_1,\ldots,g_k).
\]
The word is \emph{universal on \(G\)} if this map is surjective. In particular, for nonzero integers \(r,s\), the word \(x^ry^s\) is universal on \(G\) precisely when
\[
 \forall g\in G\ \exists x,y\in G:\quad g=x^ry^s.
\]
For a positive integer \(n\), let \(\mathbb S_n\) be the symmetric group on \(n\) letters and \(\mathbb A_n\) its alternating subgroup.

The supplied source defines
\[
 m(r,s)=
 \begin{cases}
 \displaystyle\prod_{\substack{p\text{ prime}\\p\mid rs}}p,
   & r\notin\{-1,1\}\text{ and }s\notin\{-1,1\},\\
 1,&r\in\{-1,1\}\text{ and }s\in\{-1,1\}.
 \end{cases}
\]
As written, this definition leaves \(m(r,s)\) undefined when exactly one exponent belongs to \(\{-1,1\}\). No value is assigned here in those cases.

\textbf{Source-reported background.} The record attributes to J. L. Brenner, R. J. Evans, and D. M. Silberger the following bound: for every real \(c>8/5\), there exists \(N_0=N_0(c)\) such that \(x^ry^s\), with \(rs\ne0\), is universal on \(\mathbb A_n\) whenever
\[
 n\ge\max\{N_0(c),\ c\log m(r,s)\}.
\]
The supplied examples are \(N_0(5/2)=5\) and \(N_0(2)=29\). These are reported claims from the source; the displayed formulation inherits its incomplete definition of \(m\) and its unspecified logarithm base.

\textbf{Requested statement (Yu. I. Merzlyakov, 1986).} For nonzero integer exponents \(r,s\) with at least one of \(r,s\) odd, find an analogous bound for the degree \(n\) ensuring that \(x^ry^s\) is universal on \(\mathbb S_n\). Thus the sought degree condition must imply
\[
 \forall g\in\mathbb S_n\ \exists x,y\in\mathbb S_n:\quad g=x^ry^s.
\]
The comparison with the alternating-group bound suggests seeking a logarithmic sufficient condition of the form
\[
 n\ge\max\{N,\ C\log m(r,s)\},
\]
with constants independent of \(r,s\); however, the source does not specify the desired admissible constants, whether the threshold \(c>8/5\) is also intended for \(\mathbb S_n\), or a complete definition of \(m\). This displayed form records the suggested target rather than an asserted resolution of those ambiguities.

The record attributes to M. Droste the previously known sufficient condition
\[
 n\ge\max\{6,\ 4m(r,s)-4\}
\]
under the same parity hypothesis. Its phrase that only this cruder bound is known is a source-specific historical qualification, not a new assessment here. The appended curation also reports an explicitly unverified AI claim of a logarithmic bound; it does not establish a solution.
\subsection{Short English statement}
Merzlyakov asks for a logarithmic degree bound ensuring that every permutation in \(\mathbb S_n\) can be written as \(x^ry^s\), for nonzero integer exponents with at least one odd. The supplied definition of \(m(r,s)\) omits mixed cases involving exactly one exponent equal to \(\pm1\), and the intended constants are unspecified. The registry labels the problem open while recording an unverified AI solution claim.
\subsection{Sources}
\begin{itemize}
\item Local UnsolvedMath registry, permanent ID \(700647\), exact source code \texttt{catalogue entry}, title \emph{Kourovka Notebook Problem 10.32}. The supplied record attributes the problem to Yu. I. Merzlyakov, proposed in 1986.
\item \emph{The Kourovka Notebook (Unsolved Problems in Group Theory)}, 10th issue (1986), Problem 10.32; also the supplied 21st issue edition (2026), \texttt{arXiv:1401.0300v48}. The collection description names E. I. Khukhro and V. D. Mazurov as editors.
\item J. L. Brenner, R. J. Evans, and D. M. Silberger, \emph{Proc. Amer. Math. Soc.} \textbf{96}, no. 1 (1986), 23--28. Supplied attribution for the alternating-group bound; no article title is supplied.
\item M. Droste, \emph{Proc. Amer. Math. Soc.} \textbf{96}, no. 1 (1986), 18--22. Supplied attribution for the symmetric-group bound; no article title is supplied.
\item M. Kinyon and J. Urban, \emph{Preprint} (2026), \url{https://kourovkanotebookorg.wordpress.com/wp-content/uploads/2026/09/kourovka-mathematics.pdf}. The supplied background and status review explicitly label the assertion that a logarithmic bound was found as an ``unverified AI claim'' or ``unverified AI solution.'' The review says that the October 2026 21st issue marks it with \(\#\).
\item Supplied status review checked on \texttt{2026-10-16}: canonical status \texttt{open}, evidence category \texttt{unrefereed\_claim}. This is supplied dated curation, not an independently verified literature assessment. No cited source was visited or independently checked for this draft.
\end{itemize}
\end{document}
